
\chapter{Arquitetura do Sistema}
\label{chap:arquitetura}

% Apresente uma visão geral da arquitetura do DoseCerta.
% Explique como os componentes físicos e de software
% estão organizados e como se relacionam.

\section{Visão geral da arquitetura}
\label{sec:visao_geral_arquitetura}

% Descreva os principais subsistemas do DoseCerta,
% suas responsabilidades e a divisão do sistema.

% Se houver um diagrama geral da arquitetura:
%
% \begin{figure}[htbp]
%     \centering
%     \includegraphics[width=\textwidth]{Figuras/diagrama_arquitetura}
%     \caption{Arquitetura geral do DoseCerta}
%     \label{fig:arquitetura_geral}
% \end{figure}

\section{Organização dos subsistemas}
\label{sec:subsistemas}

% Apresente os subsistemas que compõem a solução
% e explique a função de cada um.

\subsection{Subsistema mecânico}
\label{subsec:subsistema_mecanico}

% Descreva a estrutura física e os mecanismos
% responsáveis pelo armazenamento e pela dispensação.

\subsection{Subsistema eletrônico}
\label{subsec:subsistema_eletronico}

% Apresente os componentes eletrônicos,
% suas funções e as conexões relevantes.

\subsection{Subsistema embarcado}
\label{subsec:subsistema_embarcado}

% Explique o papel do sistema embarcado
% no controle e na coordenação do dispositivo.

\subsection{Backend e armazenamento de dados}
\label{subsec:backend_arquitetura}

% Descreva os serviços de backend, o armazenamento
% de dados e suas responsabilidades.

\subsection{Aplicativo e interface do usuário}
\label{subsec:aplicativo_arquitetura}

% Explique como o usuário interage com o sistema
% e quais informações são apresentadas pela interface.

\subsection{Visão computacional e inteligência artificial}
\label{subsec:ia_arquitetura}

% Descreva os módulos de visão computacional e IA,
% caso estejam integrados à arquitetura.
% Explique suas entradas, saídas e responsabilidades.

\section{Comunicação e integração entre componentes}
\label{sec:comunicacao_integracao}

% Descreva como os subsistemas trocam informações.
% Identifique interfaces, protocolos e fluxos de dados
% efetivamente utilizados no projeto.

\section{Fluxo de funcionamento do sistema}
\label{sec:fluxo_funcionamento}

% Descreva, em sequência, o funcionamento do DoseCerta
% desde a interação do usuário até a resposta do sistema.

% Se houver um diagrama de sequência ou fluxo:
%
% \begin{figure}[htbp]
%     \centering
%     \includegraphics[width=\textwidth]{Figuras/fluxo_funcionamento}
%     \caption{Fluxo de funcionamento do DoseCerta}
%     \label{fig:fluxo_funcionamento}
% \end{figure}

\section{Decisões arquiteturais}
\label{sec:decisoes_arquiteturais}

% Registre as principais decisões de arquitetura,
% as alternativas consideradas e as justificativas
% para as escolhas adotadas.
